\ifdefined\iscombined
	\lecturetitle{Petrinets}
	\setcounter{section}{4}
\else
\documentclass{article}
\usepackage{etoolbox}
\usepackage{enumitem}
\usepackage{amsthm}
\usepackage{comment}
\usepackage{amsmath}
\usepackage{amsfonts}
\usepackage{sectsty}
\usepackage{hyperref}
\usepackage{fancyhdr}
\usepackage[top=1in,bottom=1in,left=1in,right=1in]{geometry}
\usepackage{float}
\usepackage{graphicx}
\usepackage{tabularx}
\usepackage{mathtools}

\title{Lecture 4 - Petrinets}
\author{Matthew Farah, \href{mailto:farahm11@mcmaster.ca}{$\langle$farahm11@mcmaster.ca$\rangle$}, 400468251}
\date{\today}

\hypersetup{
        colorlinks=true,
        linkcolor=blue,
        filecolor=magenta,
        urlcolor=blue,
        citecolor=blue,
	anchorcolor=blue
}

\fancyhead{}
\fancyhead[L]{Matthew Farah, \href{mailto:farahm11@mcmaster.ca}{$\langle$farahm11@mcmaster.ca$\rangle$}, 400468251}
\fancyhead[R]{Lecture 4 - Petrinets}

\newcounter{example}
\newcounter{partctr}[example]

\newenvironment{example}{
	\refstepcounter{example}
	\setcounter{partctr}{0}
	\vspace{1em}
	\large\noindent\textbf{Ex.}\vspace{0.5em}
}{\vspace{1.5em}}

\newenvironment{parts}{
	\begin{enumerate}[
		label=\textbf{(\alph*)},
		leftmargin=3.3em,
		itemsep=0.4em,
		before=\large
	]
	\setcounter{enumi}{\value{partctr}}
}{
	\setcounter{partctr}{\value{enumi}}
	\end{enumerate}
}

\newenvironment{solution}{
	\vspace{0.3cm}\par\noindent\large\textbf{Solution:}\par\vspace{0.3em}
}{\vspace{0.8em}}

\renewcommand{\thesubsection}{\thesection.\alph{subsection}}
\renewcommand{\thesubsubsection}{\thesection.\alph{subsection}.\Roman{subsubsection}}
\makeatletter

\sectionfont{\fontsize{12}{12}\bfseries}
\subsectionfont{\fontsize{10}{12}\normalfont}
\subsubsectionfont{\fontsize{10}{12}\normalfont}

\newcommand{\parts@seriesname}{parts@\theexample}

\begin{document}

\maketitle
\pagestyle{fancy}

\fi

\begin{itemize}
	\item Petrinets are used to represent concurrent systems.
	\begin{itemize}
		\item Petrinets are bipartite graphs.
		\item An alternative representation to LTSs.
	\end{itemize}
	\item There exist different kinds of Petrinets.
	\item An \emph{elementary} Petrinet is composed of the following elements:
	\begin{enumerate}
		\item \emph{Places}.
		\begin{itemize}
			\item Places represent a resource (essentially a state) that the system can be in.
			\item Depicted by a circle.
		\end{itemize}
		\item \emph{Transitions}.
		\begin{itemize}
			\item Transitions represent an action that a system performs.
			\item Depicted by a rectangle.
		\end{itemize}
		\item \emph{Arcs}.
		\begin{itemize}
			\item Arcs represent directed connections between places and transitions.
			\begin{itemize}
				\item Since Petrinets are bipartitie, arcs cannot connect places to places or transitions to transitions.
			\end{itemize}
			\item Depicted by an arrow
		\end{itemize}
		\item \emph{Tokens}.
		\begin{itemize}
			\item Tokens represent the current state of the system.
			\item Tokens exist within places and move between places according to the arcs connecting places to transitions.
			\item A petrinet can have multiple tokens simultaneously.
			\item The current position of tokens in the system is depicted by a dot in the place the token resides in.
		\end{itemize}
	\end{enumerate}
	\item Relative to a given transition:
	\begin{itemize}
		\item An \emph{input} place is a place which connects to that transition.
		\item An \emph{output} place is a place which that transition connects to.
	\end{itemize}
	\item For a transition to fire, all input places to that transition must have a token inside it.
	\item Multiple transitions satisfying this condition can fire at this same time.
	\begin{itemize}
		\item This is how Petrinets model the simultaneous execution of an action.
	\end{itemize}
	\item When a transition fires, all output places are populated with a token and all input places are depopulated of their tokens.
	\item A given place can be both an input and output place of a transition.
	\item A \emph{deadlock} describes a situation where tokens are arranged in such a way that there exists no valid transition for any tokens to take.
	\item The \emph{firing sequence} of an elementary Petrinet are sequences of valid transitions that can be taken based on the initial location of tokens.
	\item The Petrinet representation of an LTS representing a concurrent system can be attained by representing the processes in parallel composition as Petrinets and ``glueing'' together shared transitions.
	\item \emph{Labelled elementary Petrinets} are essentially the same as elementary Petrinets, with the disintion being that if two transitions can be fired simultaenously, they must have distinct labels.
	\item \emph{Reachability graphs} are finite state machines used to represent the behaviour of Petrinets in LTS form.
	\begin{itemize}
		\item For a given FSP, the reachability graph of the FSP's Petrinet representation is equivalent to the FSPs LTS representation.
		\item Reachability graphs may allow for actions to be executed simultaneously.
	\end{itemize}
\end{itemize}

\begin{example}
	Represent the following FSPs:

	\begin{align*}
		\mathtt{MAKER} &= \mathtt{make} \rightarrow \mathtt{ready} \rightarrow \mathtt{MAKER} \\
		\mathtt{USER} &= \mathtt{ready} \rightarrow \mathtt{user} \rightarrow \mathtt{USER} \\
		||\mathtt{MAKER\_USER} &= \mathtt{MAKER} \; || \; \mathtt{USER}
	\end{align*}
\end{example}

\begin{parts}
	\item As LTSs.
\end{parts}

\begin{solution}
	%TODO: Fix my adding initial states.
	The LTSs for these FSPs are given by:

	\begin{figure}[H]
		\centering
		\includegraphics[width=\linewidth/2]{~/Desktop/University/Courses/SFWRENG-3BB4-Coursework/Lecture-Notes/Lecture-4/imgs/Maker-LTS.png}
		\caption{LTS representation of \texttt{MAKER}.}
	\end{figure}
	\begin{figure}[H]
		\centering
		\includegraphics[width=\linewidth/2]{~/Desktop/University/Courses/SFWRENG-3BB4-Coursework/Lecture-Notes/Lecture-4/imgs/User-LTS.png}
		\caption{LTS representation of \texttt{USER}.}
	\end{figure}
	\begin{figure}[H]
		\centering
		\includegraphics[width=\linewidth/2]{~/Desktop/University/Courses/SFWRENG-3BB4-Coursework/Lecture-Notes/Lecture-4/imgs/Maker_User-LTS.png}
		\caption{LTS representation of \texttt{MAKER\_USER}.}
	\end{figure}
\end{solution}

\begin{parts}
	\item As Petrinets.
\end{parts}

\begin{solution}
	The Petrinets for these FSPs are given by:

	\begin{figure}[H]
		\centering
		\includegraphics[width=\linewidth/2]{~/Desktop/University/Courses/SFWRENG-3BB4-Coursework/Lecture-Notes/Lecture-4/imgs/Maker-Petrinet.png}
		\caption{Petrinet representation of \texttt{MAKER}.}
	\end{figure}

	\begin{figure}[H]
		\centering
		\includegraphics[width=\linewidth/2]{~/Desktop/University/Courses/SFWRENG-3BB4-Coursework/Lecture-Notes/Lecture-4/imgs/User-Petrinet.png}
		\caption{Petrinet representation of \texttt{USER}.}
	\end{figure}

	By glueing these two Petrinets by the common transition ``\texttt{ready}'', we can easily represent the Petrinet for \texttt{MAKER\_USER} as:

	\begin{figure}[H]
		\centering
		\includegraphics[width=\linewidth/2]{~/Desktop/University/Courses/SFWRENG-3BB4-Coursework/Lecture-Notes/Lecture-4/imgs/Maker_User-Petrinet.png}
		\caption{Petrinet representation of \texttt{MAKER\_USER}.}
	\end{figure}
\end{solution}

\begin{example}
	Construct the reachability graph for the Petrinet representation of the \texttt{MAKER\_USER} FSP.
\end{example}

\begin{solution}
	The reachability graph for the Petrinet representation of \texttt{MAKER\_USER} is given by:

	\begin{figure}[H]
		\centering
		\includegraphics[width=\linewidth]{~/Desktop/University/Courses/SFWRENG-3BB4-Coursework/Lecture-Notes/Lecture-4/imgs/Reachability-Graph.png}
		\caption{Reachability graph of \texttt{MAKER\_USER} Petrinet.}
	\end{figure}
\end{solution}

\ifdefined\iscombined
\newpage
\else
  \end{document}
\fi
